% D2T-RNA v7 — Main manuscript (theoretical/computational RNA methods + statistical experimental design)
% Compile with: tectonic manuscript.tex   (SVGs pre-converted to PDF; see build_provenance)
\documentclass[11pt]{article}
\usepackage[margin=1in]{geometry}
\usepackage{amssymb,amsmath,amsthm}
\usepackage{booktabs}
\usepackage{graphicx}
\usepackage{natbib}
\usepackage{hyperref}
\usepackage{enumitem}

\newtheorem{theorem}{Theorem}
\newtheorem{corollary}{Corollary}
\newtheorem{lemma}{Lemma}
\newtheorem{proposition}{Proposition}
\newtheorem{definition}{Definition}
\newtheorem{assumption}{Assumption}
\newtheorem{remark}{Remark}

\title{Certified Collision-or-Separation Design for Finite RNA State Discrimination\\[0.5em]\large Theoretical/computational RNA methods + statistical experimental design}
\author{D2T-RNA Project}
\date{}

% SVG support: compile with --shell-escape for \includesvg, or pre-convert SVG to PDF.
\usepackage{svg}

\begin{document}
\maketitle

\begin{abstract}
In a finite registered RNA state/action/observation model, the target and rival parameter classes can
produce identical passive marginal laws, so deciding which class generated a fixed observed catalog
requires a certified and checkable design. Classical marginal-fiber, Test-Cover, and fixed-horizon-testing
arguments provide the component tools (fiber connectivity, test cover, sample-complexity scalings,
costed design), but none by itself binds them into a single pre-registered finite RNA
state/action/observation contract that emits a replayable collision-or-separation certificate over the
complete cross-class difference set together with decision-theoretic and costed-design consequences.
D2T-RNA provides, within a pre-registered finite RNA model,
exact and independently checkable collision-or-separation certificates over the complete difference set,
converts robust action-level separation into fixed-horizon finite-sample decision and budget bounds, and
returns a costed integer design with a design-class no-go certificate. All theoretical certificates are
model-conditional and are cross-checked by rational primal/dual and independent-checker evidence against
exact oracle enumeration over 11 synthetic micro-cases and 8 executed baselines. Public RNA data are used
only as a fail-closed retrospective evidence audit within fixed realized datasets, so the contribution is
a certified design method whose scope is bounded by the registered finite model.
\end{abstract}

\section{Introduction}\label{sec:intro}
\paragraph{Field problem.}
Designing a probe panel to discriminate between two candidate models of a finite RNA system reduces,
in the simplest setting, to deciding which of two parameter classes generated an observed catalog. When
the passive observation is information-lossy, the two classes can be \emph{observationally equivalent}
under the marginal map, and no passive experiment can separate them.

\paragraph{What existing methods already solve.}
Markov-basis and marginal-fiber theory characterize when two contingency tables lie in the same fiber
\citep{diaconis1998markov,deloera1995grobner}; Test-Cover minimizes a set of tests that separate
instances \citep{moret1991test}; fixed-horizon active testing and sample-complexity analysis give
scaling bounds for sequential discrimination \citep{chernoff1959sequential,hoeffding1963probability};
costed information-source selection and optimal experimental design optimize an allocation under cost
\citep{pukelsheim2006optimal}; and RNA chemical-mapping heuristics (mutate-and-map, LM2R-style) guide
probe selection from experimental data \citep{kladwang2011mutate,cordero2012importance,cheng2015consistent}.

\paragraph{Specific identifiability/design gap.}
These classical tools tell us \emph{when} separation is possible in generic settings, but not how to
produce a \emph{replayable, checkable collision-or-separation certificate} over a specific finite
registered model, and not what finite-sample decision/budget consequence follows from a certified
separation. The complete cross-class difference set and the registered nuisance coupling are not
articulated in the classical treatments.

\paragraph{Why generic cycle or Test-Cover arguments are insufficient.}
Generic alternating-cycle or Test-Cover arguments characterize fiber connectivity or test-cover
existence, but they do not (i) pin the certificate to a complete cross-class difference set $D$ rather
than a hand-picked cycle subset, (ii) make the nuisance coupling explicit and pre-registered, or
(iii) yield a directly auditable costed no-go certificate.

\paragraph{Research question.}
Within a finite registered model, can one certify collision-or-separation over the complete difference
set, convert separation into fixed-horizon finite-sample decision/budget bounds, and return either a
costed integer design or a design-class no-go certificate?

\paragraph{One-sentence core contribution.}
Within a pre-registered finite RNA state/action/observation model, D2T-RNA certifies exact
collision-or-separation over the complete cross-class difference set, converts robust action-level
separation into fixed-horizon finite-sample decision/budget bounds, and returns a costed integer design
with a design-class no-go certificate, using rational primal/dual and independent-checker evidence.

\paragraph{Evidence and exact scope.}
The theorem certificates are model-conditional and synthetic, cross-checked by exact oracle enumeration
over 11 micro-cases and 8 executed baselines. Public RNA data are used only as a fail-closed retrospective
evidence audit within fixed realized datasets.

\paragraph{Claim boundary.}
We make no claim of universal RNA identifiability, population generalization, independent-library
validation, prospective/blinded validation, or future wet-lab cost saving. The contribution is a
certified design method whose scope is the registered finite model.

\section{Related work}\label{sec:related}
\subsection{Markov-basis and marginal-fiber identifiability}
Diaconis--Sturmfels fiber theory and the Gr\"obner-basis treatment of the second hypersimplex describe
when two tables lie in the same marginal fiber and which moves connect them
\citep{diaconis1998markov,deloera1995grobner}. We use these as building blocks to articulate the
complete cross-class difference set; we do not claim a new connectivity theorem.

\subsection{Controlled sensing and fixed-horizon testing}
Fixed-horizon active hypothesis testing and minimax sample-complexity results
\citep{chernoff1959sequential,hoeffding1963probability} give scaling bounds that we bind to a registered
pair catalog and an explicit abstention rule, rather than repurposing them as a new testing theorem.

\subsection{Costed experimental design}
T-optimal and costed information-source selection \citep{pukelsheim2006optimal} motivate our costed
integer design. Our addition is an LP dual burden lower bound, an integrality gap, and a design-class
no-go certificate over the registered catalog; the LP/test-cover formulation itself is inherited.

\subsection{RNA chemical-mapping and retrospective evidence boundaries}
SHAPE/DMS mutate-and-map and LM2R-style heuristics \citep{kladwang2011mutate,cordero2012importance,
cheng2015consistent} guide probe selection. Unregistered continuous channels are excluded from all
categorical claims; public RNA data are confined to a fail-closed retrospective audit.

\section{Registered finite RNA model}\label{sec:model}
\begin{definition}[Finite registered RNA model]\label{def:model}
A \emph{finite registered RNA state/action/observation model} is a tuple
$(\Theta,A,\mathcal{Y},M,\{B_u\}_{u\in A},\mathcal{C},h,\mathcal{R})$ where:
\begin{itemize}[nosep]
  \item $\Theta=\Theta_0\sqcup\Theta_1$ is a finite set of latent states, partitioned into a
        \emph{target} class $\Theta_0$ and a \emph{rival} class $\Theta_1$; the two classes are
        \emph{disjoint} ($\Theta_0\cap\Theta_1=\emptyset$), so no pair in
        Definition~\ref{def:D} is the trivial diagonal pair $\theta_0=\theta_1$;
  \item $A$ is a finite set of \emph{actions} (probes/perturbations);
  \item $\mathcal{Y}$ is a finite categorical \emph{observation alphabet};
  \item $M:\Theta\to\Delta(\mathcal{Y})$ is the \emph{passive marginal map};
  \item for each $u\in A$, $B_u:\Theta\to\Delta(\mathcal{Y})$ is the \emph{action-level observation law};
  \item $\mathcal{C}$ is a registered \emph{nuisance coupling} (CARTESIAN or
        EQUAL\_REALIZED\_VALUE) that authorizes the action-level product law;
  \item $h\in\mathbb{Z}_{\ge 0}$ is a fixed \emph{non-adaptive horizon};
  \item $\mathcal{R}$ is an \emph{abstention rule} (decision-or-abstain).
\end{itemize}
\end{definition}

\begin{definition}[Complete cross-class difference set $D$]\label{def:D}
The \emph{complete cross-class difference set} is
$D=\Theta_0\times\Theta_1$, the set of all target/rival pairs. It is the full cross product, not a
hand-picked cycle subset. Because $\Theta_0$ and $\Theta_1$ are finite and disjoint
(Definition~\ref{def:model}), $D$ is finite and contains no trivial diagonal pair $(\theta,\theta)$.
\end{definition}

\begin{remark}[Relation to the contract difference set]\label{rem:D-relation}
This paper's $D=\Theta_0\times\Theta_1$ is the finite-catalog representation of the contract's
cross-class difference set: each pair $(\theta_0,\theta_1)\in D$ indexes the difference of the two
latent state distributions, and certifying separation over the \emph{full} cross product (rather than
only the passive collision fiber, i.e.\ those pairs whose passive marginal laws coincide) is
a conservative, stronger certificate. It is the full cross product, not a hand-picked subset.
\end{remark}

\begin{remark}[Finite catalog, attainment, and triviality]\label{rem:attain}
The paper restricts the model to \emph{finite} state classes (finite catalog), so $D$ is finite and the
minimum in Definition~\ref{def:gamma} is attained. Consequently $\gamma(S)=0$ always yields a
\emph{non-trivial, attained} exact collision witness $(\theta_0,\theta_1)\in D$ with
$\theta_0\neq\theta_1$ (disjointness excludes the diagonal). The contract's general framework also
allows \emph{rational polyhedral} classes, in which case $D$ is a difference set of state
distributions and the infimum may not be attained; there the equivalence degrades to
necessary-only / sufficient-only forms and the infimum-only obstruction boundary (contract 5.2), which
is out of scope here. We therefore do \emph{not} claim an attained witness when the infimum is only
approached.
\end{remark}

\begin{definition}[Action panel and induced law]\label{def:panel}
For a panel $S\subseteq A$ of actions, the \emph{induced observation law} of state $\theta$ is the
product law $P_S(\theta)=\bigotimes_{u\in S}B_u(\theta)$, where the product is authorized by the
registered coupling $\mathcal{C}$.
\end{definition}

\begin{remark}[Passive marginal, induced law, and the action map]\label{rem:passive}
$M$ is the \emph{passive} (no-action) observation law and defines the collision fiber: a pair
$(\theta_0,\theta_1)$ is passively collision-equivalent exactly when the two induced passive laws
coincide. The action map $B_u$ induces, on the categorical alphabet $\mathcal{Y}$, the atomic
observation law used to \emph{separate} the pair; the complete \emph{categorical observation law} for a
panel $S$ is the product law $P_S(\theta)=\bigotimes_{u\in S}B_u(\theta)$ over
$\mathcal{Y}^{|S|}$,
authorized by the registered coupling $\mathcal{C}$. Thus $\gamma(S)$ measures, in
\emph{observation-probability} terms, how well the action maps separate every target/rival pair in $D$.
\end{remark}

\begin{definition}[Separation functional $\gamma(S)$]\label{def:gamma}
$\gamma(S)=\min_{(\theta_0,\theta_1)\in D}\mathrm{TV}\bigl(P_S(\theta_0),P_S(\theta_1)\bigr)$,
the minimum total-variation separation of the induced laws over the complete difference set. Explicitly,
$\gamma(S)$ is a norm on pairs of \emph{categorical observation-probability} laws over
$\mathcal{Y}^{|S|}$ (the product law $P_S$ of Definition~\ref{def:panel}), \emph{not} a metric on the
state space $\Theta$. Because $D$ is finite and the laws are finite categorical laws, the minimum is
attained (Remark~\ref{rem:attain}).
\end{definition}

\begin{assumption}[Registered structure]\label{ass:reg}
The model is pre-registered: the state classes, action set, alphabet, marginal and action maps, coupling
$\mathcal{C}$, horizon $h$, and abstention rule $\mathcal{R}$ are fixed before any data are used.
\end{assumption}

\begin{assumption}[Complete action-level likelihoods]\label{ass:likelihood}
For every pair in $D$ and every action $u\in A$, the action-level likelihoods $B_u(\theta_0)(y)$ and
$B_u(\theta_1)(y)$ are known and positive, so the product law is fully specified.
\end{assumption}

\paragraph{Explicitly not covered.}
Continuous infinite-dimensional nuisance, adaptive actions, randomized design, sequential stopping,
arbitrary high-dimensional contingency tables, unregistered continuous SHAPE/DMS channels, unknown
action-induced state change, and unmodeled third states.

\section{Exact geometry (T2b)}\label{sec:t2b}
\begin{theorem}[Exact collision-or-separation, T2b]\label{thm:t2b}
Let the model of Definition~\ref{def:model} satisfy Assumptions~\ref{ass:reg}--\ref{ass:likelihood}.
For a panel $S\subseteq A$:
\begin{enumerate}[nosep]
  \item $\gamma(S)=0$ \emph{if and only if} there exists $(\theta_0,\theta_1)\in D$ such that
        $P_S(\theta_0)=P_S(\theta_1)$ (an \emph{exact collision witness});
  \item $\gamma(S)>0$ \emph{if and only if} $P_S(\theta_0)\neq P_S(\theta_1)$ for every
        $(\theta_0,\theta_1)\in D$ (a \emph{separation certificate}).
\end{enumerate}
Moreover, when $\gamma(S)>0$, the pair attaining the minimum yields a rational primal/dual certificate
and an independent-checker verification.
\end{theorem}

\begin{proof}[Proof sketch]
For finite $D$, the map $(\theta_0,\theta_1)\mapsto\mathrm{TV}(P_S(\theta_0),P_S(\theta_1))$ takes
finitely many values, so its minimum is attained (Definition~\ref{def:gamma}). If the attained minimum is
$0$, some pair has $\mathrm{TV}=0$, i.e.\ identical product laws, an exact collision witness, giving the
forward direction of (i); the reverse direction is immediate. If the minimum is $>0$, every pair has
positive TV, i.e.\ distinct product laws, a separation certificate; the reverse direction is immediate.
The rationality follows because all laws are finite rational categorical laws and the certificate is
witnessed by an explicit pair; an independent checker re-derives the collision or separation from first
principles (\texttt{verify.py verify\_collision}/\texttt{verify\_separation}) and the LP
strong-duality + primal feasibility check supplies the rational primal/dual pair
(\texttt{lp.py strong-duality}). The full proof is in the supplementary.
\end{proof}

\begin{remark}[Direction semantics]\label{rem:direction}
For finite catalogs the infimum in $\gamma(S)$ is attained, so the theorem is \emph{iff}. Both
registered couplings ($\mathcal{C}=\textsc{CARTESIAN}$ and $\mathcal{C}=\textsc{EQUAL\_REALIZED\_VALUE}$)
fully specify the product law under Assumption~\ref{ass:likelihood}, so the iff holds under either
coupling. The necessary-only / sufficient-only forms are reserved for the non-atomic / infimum-only
obstruction boundary (contract 5.2), which we do not treat here (Remark~\ref{rem:attain}).
\end{remark}

\begin{remark}[What the rational certificate verifies]\label{rem:certificate}
The certificate is a statement about \emph{observation-probability separation}, not a purely geometric
($\Theta$) claim. The attaining pair $(\theta_0^\star,\theta_1^\star)\in D$ with
$\mathrm{TV}(P_S(\theta_0^\star),P_S(\theta_1^\star))=\gamma(S)$ is a witnessed element of $D$, and the
value $\gamma(S)>0$ is the categorical-law separation it certifies. The LP dual is an
\emph{intermediate} certificate for the costed-design consequence (T2d): it certifies that every
feasible design has cost at least the dual bound, not a second, independent geometric claim. The
independent checker re-derives both the witnessed collision/separation and the dual from first
principles (\texttt{verify.py} and \texttt{lp.py}).
\end{remark}

\begin{figure}[ht]
\centering
\includesvg[width=0.85\textwidth]{figures/fig1_geometry}
\caption{Exact geometry (T2b). \textbf{Question:} does the induced law map separate the complete
difference set $D$? \textbf{Result:} oracle minimax decision error per synthetic microcase
(model-conditional); colliding pairs reach error $1/2$, separated pairs reach error $0$.
\textbf{Interpretation:} exact collision-or-separation certificate consistent with the oracle.
\textbf{Boundary:} model-conditional; no population generalization.}
\label{fig:geometry}
\end{figure}

\section{Finite-sample consequence (T2c)}\label{sec:t2c}
\begin{definition}[Pair information]\label{def:info}
For a pair $(\theta_0,\theta_1)\in D$ and action $u$, the
\emph{Hellinger/Bhattacharyya information} \citep{hellinger1909neue} is
$I_u(\theta_0,\theta_1)=-\log\sum_{y\in\mathcal{Y}}\sqrt{B_u(\theta_0)(y)\,B_u(\theta_1)(y)}$.
Under $n_u$ repeated non-adaptive observations of action $u$ with the registered product law, the total
information is $I_{\mathrm{tot}}=\sum_{u\in A}n_u\,I_u(\theta_0,\theta_1)$.
\end{definition}

\begin{definition}[Decision-theoretic endpoints]\label{def:endpoints}
Fix a \emph{decision-or-abstain} rule $\mathcal{R}$ that, on each realized catalog, returns either a
declared class or \emph{inconclusive} (abstain). The registered endpoints are:
\begin{itemize}[nosep]
  \item $\alpha$: the type-I wrong-declaration bound (declaring the rival when the target generated the data);
  \item $\beta$: the type-II wrong-declaration bound (declaring the target when the rival generated the data);
  \item $\kappa_0$: the minimum correct-declaration probability under the target hypothesis (H0);
  \item $\kappa_1$: the minimum correct-declaration probability under the rival hypothesis (H1);
  \item $\rho=\mathbb{P}(\text{abstain})$: the probability the rule returns \emph{inconclusive}.
\end{itemize}
Every declaration is exactly one of correct, wrong, or abstained; abstention is neither correct nor
wrong, so $\mathbb{P}(\mathrm{correct})+\mathbb{P}(\mathrm{wrong})+\rho=1$.
\end{definition}

\begin{theorem}[Finite-sample decision bound, T2c]\label{thm:t2c}
Consider a finite registered pair catalog under a fixed-horizon deterministic non-adaptive allocation
with complete action-level likelihoods and registered product laws, with endpoints as in
Definition~\ref{def:endpoints}.
\begin{enumerate}[nosep]
  \item \emph{(Upper bound.)} For every pair, the probability of a wrong declaration satisfies
        $\mathbb{P}(\mathrm{wrong})\le \tfrac12\exp(-I_{\mathrm{tot}})$, so in particular
        $\alpha,\beta\le\tfrac12\exp(-I_{\mathrm{tot}})$; the correct-declaration probability satisfies
        $\mathbb{P}(\mathrm{correct})\ge 1-\tfrac12\exp(-I_{\mathrm{tot}})-\rho$, where $\rho$ is the
        abstention probability under rule $\mathcal{R}$.
  \item \emph{(Lower bound / structural no-go.)} To reach a correct-declaration floor
        $\kappa=\max(\kappa_0,\kappa_1)\in(1/2,1)$, the total information must satisfy
        $I_{\mathrm{tot}}\ge -\tfrac12\log\!\bigl(1-(2\kappa-1)^2\bigr)$. If the certified total
        information is strictly below this threshold, no allocation within the budget can meet the
        floor.
\end{enumerate}
\end{theorem}

\begin{proof}[Proof sketch]
For (i), the Bhattacharyya bound gives
$\mathbb{P}(\mathrm{wrong})\le\tfrac12\prod_{u\in A}\rho_u^{n_u}$, where
$\rho_u=\sum_y\sqrt{B_u(\theta_0)(y)B_u(\theta_1)(y)}=e^{-I_u}$, so the product is
$e^{-I_{\mathrm{tot}}}$. For (ii), the decision error is minimized when the two product laws are
separated by the largest achievable TV; relating the required separation to the Bhattacharyya coefficient
and rearranging the floor condition $\mathbb{P}(\mathrm{correct})\ge\kappa$ yields the stated
information threshold. The arithmetic uses exact rationals (\texttt{fractions.Fraction}) and
directed-rounding decimal intervals (\texttt{Decimal}) for the transcendental terms; the upper bound is
cross-checked against exact minimax enumeration (\texttt{decision.py exact\_minimax\_error}). The full
proof is in the supplementary.
\end{proof}

\begin{remark}[Tightness microcase]\label{rem:tight}
For alphabet size 2, $P_0=(1/4,3/4)$, $P_1=(1,0)$: the Bhattacharyya coefficient is $1/2$, the
information is $\ln2$, the exact minimax error is $\tfrac12(1/4)^n$, and the certified upper bound
$\tfrac12\exp(-nI)$ brackets the Bhattacharyya-ideal $\tfrac12(1/2)^n$, i.e.\ the finite-sample bound
dominates the exact minimax error and is conservative by the intended margin.
\end{remark}

\begin{remark}[Budget and sample complexity]\label{rem:budget}
To meet a correct-declaration floor $\kappa=\max(\kappa_0,\kappa_1)$ with abstention $\rho$, T2c~(ii)
requires $I_{\mathrm{tot}}\ge -\tfrac12\log\!\bigl(1-(2\kappa-1)^2\bigr)$. Expressed as a budget, the
minimal non-adaptive allocation is the integer solution of
$\min_{n\in\mathbb{Z}_{\ge0}^{A}}\sum_{u\in A}c_u n_u$ subject to
$\sum_{u\in A}n_u I_{uw}\ge -\tfrac12\log\!\bigl(1-(2\kappa_w-1)^2\bigr)$ for every active pair $w$,
which is exactly the T2d program (Section~\ref{sec:t2d}) with $\tau_w$ set to that per-pair threshold.
Thus T2c supplies the per-pair information/budget target and T2d turns it into a costed design or a
design-class no-go. The bounds are stated for the finite registered pair catalog only; no
composite-continuous covering or real-data calibration guarantee is made.
\end{remark}

\begin{figure}[ht]
\centering
\includesvg[width=0.85\textwidth]{figures/fig2_finiteness}
\caption{Finite-sample consequence (T2c). \textbf{Question:} how does the certified decision bound
compare with exact oracle enumeration? \textbf{Result:} baseline minimax decision error by exact oracle
enumeration for a representative separated case (model-conditional). \textbf{Interpretation:} the
finite-sample decision consequence cross-checks the exact enumeration. \textbf{Boundary:} no
composite-continuous covering theorem, no arbitrary continuous nuisance guarantee, no real-data
calibration.}
\label{fig:finiteness}
\end{figure}

\section{Costed design and no-go (T2d)}\label{sec:t2d}
\begin{definition}[Costed design problem]\label{def:costed}
Given unit costs $c_u\ge 0$ and requirement thresholds $\tau_w\ge 0$ over active pairs $w\in D'$, the
\emph{costed design problem} is
\[
\min_{n\in\mathbb{Z}_{\ge 0}^{A}}\ \sum_{u\in A}c_u n_u
\quad\text{s.t.}\quad \sum_{u\in A}n_u I_{uw}\ge \tau_w\ \ \forall w\in D',
\]
where $I_{uw}$ is the certified pair-information interval for pair $w$ under action $u$.
\end{definition}

\begin{definition}[LP relaxation and dual]\label{def:lp}
The \emph{LP relaxation} relaxes $n_u\ge 0$ to real. Its \emph{dual} is
$\max_{y\ge 0}\sum_{w\in D'}\tau_w y_w$ subject to $\sum_{w\in D'}I_{uw}\,y_w\le c_u$ for all $u\in A$.
For any dual-feasible $y$, $\sum_w\tau_w y_w$ is a \emph{dual burden lower bound} on the cost of every
feasible design.
\end{definition}

\begin{theorem}[Costed design and no-go, T2d]\label{thm:t2d}
Let the model satisfy Assumptions~\ref{ass:reg}--\ref{ass:likelihood}.
\begin{enumerate}[nosep]
  \item \emph{(Achievable upper.)} If an integer point $n^\star$ satisfies
        $\sum_u n^\star_u I_{uw}^{\mathrm{lo}}\ge \tau_w$ for all $w\in D'$, then $c^T n^\star$ is an
        achievable cost (upper bound) and its feasibility is explicitly checkable.
  \item \emph{(No-go lower.)} If a dual-feasible $y$ gives $\tau^T y>\mathrm{budget}$, then no feasible
        design within budget exists (a \emph{design-class no-go certificate}).
  \item \emph{(Integrality gap.)}
        $\mathrm{gap}=\bigl(\text{upper cost}-\text{lower bound}\bigr)/\text{lower bound}$, computed
        conservatively with certified interval information.
\end{enumerate}
\end{theorem}

\begin{proof}[Proof sketch]
(i) is by construction: any integer $n^\star$ satisfying the coverage constraints is feasible and has
cost $c^Tn^\star$; feasibility is re-checked from first principles
(\texttt{costed\_verify.check\_integer\_design\_feasible}). (ii) follows by weak duality: every feasible
design satisfies $\sum_u c_u n_u\ge\sum_w\tau_w y_w$ for every dual-feasible $y$ (using the interval
upper bound on $I_{uw}$ keeps the LP a relaxation), so a dual bound exceeding the budget rules out all
feasible designs. (iii) is direct arithmetic. The full proof is in the supplementary.
\end{proof}

\begin{remark}[Exact microcases]\label{rem:costed}
\emph{Single-pair integral:} $\min n$ with $n\ge 3$, cost $2$ gives IP$=$LP$=6$, gap $0$.
\emph{Fractional cover $3\times3$:} $A=[[1,0,1],[1,1,0],[0,1,1]]$, $\tau=(1,1,1)$, $c=(1,1,1)$ gives
LP$=3/2$, IP$=2$, gap $=1/3$. \emph{Fractional requirement:} $\min n$ with $n\ge 3/2$, cost $1$ gives
LP$=3/2$, IP$=2$, gap $=1/3$. \emph{No-go demonstration:} a fractional cover with budget $7/5<3/2=$ LP
bound is NO\_GO.
\end{remark}

\begin{remark}[No-go logic and scope]\label{rem:nogo}
The no-go is a \emph{design-class} certificate: a dual bound above budget rules out \emph{every}
feasible design within budget, even though the LP relaxation may itself be feasible ---
LP-feasible does not imply integer-feasible. The certificate is not a claim that no RNA assay exists;
it certifies the registered costed design class over the finite pair catalog $D'$. The integrality-gap
values above make the ``LP-feasible but not integer-feasible'' gap concrete (gap $=1/3$ for the
fractional cases).
\end{remark}

\begin{figure}[ht]
\centering
\includesvg[width=0.85\textwidth]{figures/fig3_costed_design}
\caption{Costed design (T2d). \textbf{Question:} when is a costed design feasible, and when is it a
design-class no-go? \textbf{Result:} log-scale LP dual lower bound per microcase, with achieved integer
upper cost where feasible; infeasible cases are marked no-go. \textbf{Interpretation:} supports the
costed design and design-class no-go certificate. \textbf{Boundary:} no globally optimal RNA design, no
universal minimum-cost strategy, no wet-lab cost already saved.}
\label{fig:costed}
\end{figure}

\section{Worked numerical cases}\label{sec:worked}
We report the concrete numbers behind the three theorems, drawn from the executed
model-conditional evaluation matrix (11 micro-cases $\times$ 8 baselines; full 88-row table in the
supplementary).

\begin{table}[ht]
\centering
\caption{Worked cases (T2 integer+LP design). \textbf{Question:} do the certificates and bounds
reproduce at the level of concrete numbers? \textbf{Result:} cost, product TV, minimax error, and
correct-declaration rate per representative micro-case. \textbf{Interpretation:} exact collision
(resp.\ strict separation, cancellation, near-collision) reproduce $\gamma=0$ (resp.\ $\gamma>0$,
$\gamma=0$ via cycle cancellation, $\gamma>0$ small); the finite-sample/by-oracle comparison and the
cost/no-go case quantify the decision and design consequences. \textbf{Boundary:} model-conditional
synthetic evaluation; no method-superiority or population claim.}
\label{tab:worked}
\begin{tabular}{@{}llccccc@{}}
\toprule
Required case & micro-case & cost & TV & minmax err & correct & verdict \\
\midrule
exact collision & exact\_collision & 8 & 0 & 1/2 & 1/2 & $\gamma=0$, witness \\
strict separation & strict\_positive\_separation & 1 & 1 & 0 & 1 & $\gamma>0$, cert. \\
cancellation & cancellation\_cycled & 8 & 0 & 1/2 & 1/2 & $\gamma=0$ (cancellation) \\
finite-sample vs exact oracle & boundary & 1 & 1/2 & 1/4 & 3/4 & $8\to1$ cost, exp. bound \\
cost / no-go & alternating\_rectangle\_2x2 & 8 & 0 & 1/2 & 1/2 & no integer upper, no-go \\
\bottomrule
\end{tabular}
\end{table}

\paragraph{Finite-sample bound vs exact oracle.}
For the T2c tightness microcase ($P_0=(1/4,3/4)$, $P_1=(1,0)$, alphabet 2), the exact minimax error is
$\tfrac12(1/4)^n$ while the certified upper bound brackets $\tfrac12(1/2)^n$
(Remark~\ref{rem:tight}); the certified bound dominates the exact oracle error for every $n\ge1$.

\paragraph{Abstention boundary.}
The assumption gate (section 10) evaluates six violation profiles (action changes state, omitted third
state, unregistered nuisance, unknown dependency unit, no complete observation model, dependent
observations). Under a valid registered profile the certificate proceeds (PROCEED); under any violation
profile the decision abstains as NOT\_ESTABLISHED. No input condition is force-fit into the
theorem-required form. This is a theory boundary, not real-data validation.

\section{Synthetic evaluation and baselines}\label{sec:eval}
Eleven micro-cases and eight baselines are executed as model-conditional synthetic evaluation.
Table~\ref{tab:baselines} states the comparison variables on which every baseline is evaluated. Executed
baselines do not imply biological superiority; oracle agreement does not imply population
generalization; runtime/memory are computational evidence only.

\paragraph{Definition of metrics.}
Fix an active pair $(\theta_0,\theta_1)\in D'$ and a baseline design (a selected panel $S$ with a
decision-or-abstain rule). \emph{Product TV} is the total-variation distance between the induced
categorical observation laws, $\mathrm{TV}(P_S(\theta_0),P_S(\theta_1))$. \emph{Minimax error} is the
exact two-class minimax decision error of the oracle,
$\tfrac12\bigl(1-\mathrm{TV}(P_S(\theta_0),P_S(\theta_1))\bigr)$, which is attained by the likelihood-ratio
test under the registered product law. \emph{Correct} and \emph{abstain} are the empirical
correct-declaration probability $\mathbb{P}(\mathrm{correct})$ and abstention probability
$\rho=\mathbb{P}(\mathrm{abstain})$ of the baseline rule (Definition~\ref{def:endpoints}). Because the
oracle minimax error is an affine decreasing function of product TV, the two columns are not independent;
we report both for auditability. $\mathrm{abstain}=0$ means the rule always declares.

\begin{table}[ht]
\centering
\caption{Baseline comparison variables. \textbf{Question:} on what common variables are all eight
baselines compared? \textbf{Result:} every baseline is scored on cost, product TV, minimax error,
abstention, and correct-declaration rate; the T2 integer+LP baseline additionally reports the LP dual
lower bound, integer upper cost, and integrality gap. \textbf{Interpretation:} the comparison is
well-defined because all baselines share the same registered action/cost space. \textbf{Boundary:}
same registered action/cost space only; no cross-space or biological comparison.}
\label{tab:baselines}
\begin{tabular}{@{}llllll@{}}
\toprule
Baseline & cost & product TV & minmax err & abstain & corr. decl. \\
\midrule
exhaustive\_oracle & \checkmark & \checkmark & \checkmark & \checkmark & \checkmark \\
full\_matrix & \checkmark & \checkmark & \checkmark & \checkmark & \checkmark \\
random & \checkmark & \checkmark & \checkmark & \checkmark & \checkmark \\
greedy\_test\_cover & \checkmark & \checkmark & \checkmark & \checkmark & \checkmark \\
eig & \checkmark & \checkmark & \checkmark & \checkmark & \checkmark \\
chernoff & \checkmark & \checkmark & \checkmark & \checkmark & \checkmark \\
lm2r\_heuristic & \checkmark & \checkmark & \checkmark & \checkmark & \checkmark \\
t2\_integer\_lp & \checkmark & \checkmark & \checkmark & \checkmark & \checkmark \\
\bottomrule
\end{tabular}
\\[0.5em]
\footnotesize T2 integer+LP additionally reports LP dual lower bound, integer upper cost, and
integrality gap. Full 88-row table is in the supplementary.
\end{table}

\begin{table}[ht]
\centering
\caption{Abbreviated baseline summary (cost, oracle minimax error) for representative key cases.
The complete 88-row table is in the supplementary.}
\label{tab:abbrev}
\begin{tabular}{llccccc}
\toprule
case & method & cost & err & tv & abstain & correct \\
\midrule
exact\_collision & t2\_integer\_lp & 8 & 1/2 & 0 & 0 & 1/2 \\
boundary & t2\_integer\_lp & 1 & 1/4 & 1/2 & 0 & 3/4 \\
strict\_positive\_separation & t2\_integer\_lp & 1 & 0 & 1 & 0 & 1 \\
no\_cycle & t2\_integer\_lp & 1 & 0 & 1 & 0 & 1 \\
\bottomrule
\end{tabular}
\end{table}

\section{Retrospective evidence qualification (Task6-R)}\label{sec:retro}
This section is \emph{evidence qualification}, not validation. The three registered datasets close
fail-closed with distinct terminal outcomes; no established quantitative instance is claimed.
Table~\ref{tab:retro} states, per dataset, the registered role, the R2 outcome, what can be reported,
and what cannot be claimed.

\begin{table}[ht]
\centering
\caption{Retrospective evidence qualification. \textbf{Question:} how do the registered datasets
close, and what may/may not be reported? \textbf{Result:} add, SAM-III, and RORC close fail-closed with
distinct terminal outcomes. \textbf{Interpretation:} data are used only as fixed realized catalogs for a
fail-closed audit. \textbf{Boundary:} no categorical replay, no quantitative T2 validation, no future
cost saving, no independent validation.}
\label{tab:retro}
\begin{tabular}{@{}p{2.0cm}p{3.2cm}p{3.4cm}p{3.2cm}p{3.2cm}@{}}
\toprule
Dataset & Registered role & R2 outcome & What can be reported & What cannot be claimed \\
\midrule
add/RMDB & retrospective full-matrix compression if qualified &
  NOT\_COMPARABLE\_BY\_REGISTERED\_OBSERVATION\_MODEL &
  modality limitation; raw provenance; fail-closed audit &
  categorical replay; quantitative T2 validation; future cost saving \\
\addlinespace
SAM-III/GSE278422 & modality-transfer diagnostic &
  NOT\_COMPARABLE\_BY\_REGISTERED\_ACTION\_SPACE &
  action/modality incompatibility; diagnostic failure &
  unified benchmark; quantitative transfer \\
\addlinespace
RORC & exposed misspecification stress &
  NOT\_APPLICABLE &
  explicit terminal audit and boundary semantics &
  independent validation; unseen OOD; third-state discovery \\
\bottomrule
\end{tabular}
\end{table}

\begin{figure}[ht]
\centering
\includesvg[width=0.9\textwidth]{figures/fig4_retro}
\caption{Retrospective evidence qualification (Task6-R). \textbf{Question:} how do the three
registered datasets close? \textbf{Result:} each closes fail-closed with a distinct terminal outcome
(observation-model, action-space, applicability). \textbf{Interpretation:} data serve only as a
fail-closed audit within fixed realized datasets. \textbf{Boundary:} no established quantitative
instance is claimed.}
\label{fig:retro}
\end{figure}

\section{Misspecification and abstention}\label{sec:misspec}
Valid registered assumptions yield PROCEED; assumption violation yields NOT\_ESTABLISHED or abstain.
This is a theory boundary and credibility evidence, not real-data validation.

\begin{figure}[ht]
\centering
\includesvg[width=0.85\textwidth]{figures/fig5_misspec}
\caption{Misspecification boundary and abstention. \textbf{Question:} when does the certificate
proceed versus abstain? \textbf{Result:} valid registered assumptions yield PROCEED; any of the six
violation profiles yields abstain/NOT\_ESTABLISHED. \textbf{Interpretation:} no input condition is
force-fit into the theorem-required form (contract 10.2). \textbf{Boundary:} theory boundary and
credibility evidence, not real-data validation.}
\label{fig:misspec}
\end{figure}

\section{Discussion}\label{sec:discussion}
What is genuinely new is the replayable, checked, model-conditional certificate framework over the
complete registered difference set with finite-sample and costed design consequences. What is inherited
is the classical marginal-fiber, Test-Cover, and testing toolbox (Sections~\ref{sec:related}).
The RNA-feasible composite formulation is useful because it makes the identifiability/design question
decision-theoretic and auditable. The finite-sample consequence adds a quantitative bound tied to action
geometry. Without a new library or prospective experiment, no population-level, independent-validation,
or third-state-discovery claim can be made.

\section*{Supplementary material}
Full proofs of Theorems~\ref{thm:t2b}--\ref{thm:t2d}, the exact micro-cases, and the complete 88-row
baseline table are provided in the supplementary document.

\bibliographystyle{plainnat}
\bibliography{references}

\end{document}